\label{chap:introduction}

Recent technological advances have led manufacturing companies to furnish their machines with increasingly sophisticated sensors and interfaces. At the same time, the widespread adoption of Artificial Intelligence (AI) technologies, particularly Neural Networks (NNs), has encouraged companies to investigate how these methods could improve machine functionalities.

As a result, these technological developments have enabled the creation of more advanced machinery integrating Intelligent Decision Support Systems (IDSS)~\cite{phillips2013intelligent}, capable of assisting operators during manufacturing activities.

Although the machinery produced varies significantly in functionality and features, many manufacturing processes require interaction with a human operator and compliance with operational and safety constraints. During the initial stage of production, the operator typically configures the parameters governing the manufacturing process, while in the final stage the quality of the finished product is evaluated.

The physical processes underlying manufacturing operations directly influence the quality of the final product. However, operators often lack the time or expertise required to analyze these processes in detail and therefore rely on the machine interface to provide guidance and support for decision making.

The objective of this work is to improve the decision-making process of human operators interacting with manufacturing machines by developing Intelligent Decision Support Systems (IDSS) aimed at optimizing the tasks performed during machine configuration.

This doctoral thesis was developed in collaboration with SCM Group S.p.A., an Italian company and world leader in the production of machinery for processing a wide range of materials, including wood, plastic, glass, stone, metal, and composite materials.

The remainder of this thesis investigates two optimization problems arising in industrial machine configuration. The first, referred to as \textit{Fixture Layout Optimization}, addresses the positioning and configuration of fixtures used to support wooden workpieces during machining in a \textit{work center}. The problem is investigated by integrating exact optimization methods with heuristic search techniques to identify feasible, high-quality fixture configurations. This problem is addressed in the first part of the thesis.

The second problem, referred to as \textit{Heating Process Optimization}, concerns the determination of machine parameters required to achieve a target temperature distribution within the core of a plastic sheet in a \textit{thermoforming} machine. The problem is addressed using neural networks trained on data generated by a physics-based simulator, whose predictions are subsequently refined through numerical optimization to improve accuracy while maintaining cost-efficient process configurations. This problem is addressed in the second part of the thesis.

